% Shared preamble for the Generative AI (CS 782/882, Fall 2026) lecture notes.
% Each weekNN-notes.tex file inlines this block so it compiles on its own.
\documentclass[11pt]{article}
\usepackage[margin=1in]{geometry}
\usepackage[T1]{fontenc}
\usepackage[utf8]{inputenc}
\usepackage{lmodern}
\usepackage{amsmath,amssymb,amsthm,mathtools,bm}
\usepackage{enumitem}
\usepackage{graphicx}
\usepackage{booktabs}
\usepackage{xcolor}
\usepackage{tikz}
\usetikzlibrary{arrows.meta,positioning,calc,shapes}
\usepackage{listings}
\usepackage[colorlinks=true,linkcolor=blue!60!black,urlcolor=blue!60!black,citecolor=blue!60!black]{hyperref}

\newcommand{\E}{\mathbb{E}}
\newcommand{\Var}{\mathrm{Var}}
\newcommand{\KL}{\mathrm{KL}}
\newcommand{\given}{\,\vert\,}
\newcommand{\R}{\mathbb{R}}
\DeclareMathOperator*{\argmax}{arg\,max}
\DeclareMathOperator*{\argmin}{arg\,min}
\newtheorem{theorem}{Theorem}
\newtheorem{lemma}{Lemma}
\newtheorem{definition}{Definition}
\theoremstyle{remark}
\newtheorem*{remark}{Remark}

\lstset{basicstyle=\ttfamily\small,breaklines=true,frame=single,columns=fullflexible}

% Additional definitions for this supplementary note

\title{Supplementary Notes: Teacher-Student Distillation}
\author{Hong Qin, Old Dominion University}
\date{Fall 2026}

\begin{document}
\maketitle

\begin{quote}\small\itshape Note. The material in these supplementary notes was covered in the 16-week version of this course and is \textbf{optional reading} for the 8-week offering. It assumes the Week 6 lecture notes on transformers (self-attention, positional encoding, the transformer block, and next-token training) and is provided for students who want to see how those ideas are extended in current research systems.\end{quote}

\tableofcontents
\newpage

\section{Conceptual Summary}

Reference: G. Hinton, O. Vinyals, J. Dean, ``Distilling the Knowledge in a Neural Network,'' 2015, arXiv:\allowbreak 1503.02531.

\begin{itemize}
    \item \textbf{Goal:} Transfer a distribution $q_\phi(\cdot \mid c)$ learned by a large/accurate \emph{teacher} to a smaller/faster \emph{student} $p_\theta(\cdot \mid c)$.
    \item \textbf{Core idea:} Make the student match the teacher's \emph{soft targets} (probability distributions, geometry distributions, or denoising trajectories), optionally combined with ground-truth \emph{hard targets} and/or \emph{intermediate feature} matching.
    \item \textbf{Why it works:} Soft targets encode ``dark knowledge'' (relative similarities among outcomes) that hard labels don't convey; this often improves generalization and compresses large models effectively.
\end{itemize}

\section{Unified Objective}
Let
\begin{itemize}[nosep]
    \item $c$ = context (input $x$, history $y_{<t}$, conditioning variables),
    \item $o$ = predicted object (class, token, structure, noise, score, \ldots),
    \item $q_\phi(o \mid c; T)$ = teacher distribution with temperature $T>1$,
    \item $p_\theta(o \mid c; T)$ = student distribution.
\end{itemize}

We minimize:
\[
\min_{\theta} \;
\E_{c \sim \mathcal{D}} \Big[
    \alpha \, \ell_{\text{hard}}(c;\theta)
    + \beta \, T^{2} \, D\!\big(q_\phi(\cdot \mid c;T)\,\|\,p_\theta(\cdot \mid c;T)\big)
    + \gamma \, R_{\text{feat}}(c;\theta,\phi)
    + \delta \, R_{\text{traj}}(c;\theta,\phi)
\Big].
\]

\textbf{Components:}
\begin{itemize}
    \item $\ell_{\text{hard}}$: ground-truth loss (CE/NLL/MSE). Set $\alpha=0$ for pure distillation.
    \item $D$: divergence (KL, reverse-KL, JS, Wasserstein, energy distance).
    \item $R_{\text{feat}}$: hidden state matching, e.g. $\|h_T(c)-h_S(c)\|^2$.
    \item $R_{\text{traj}}$: trajectory matching (diffusion timesteps, autoregressive rollouts).
\end{itemize}

\section{Canonical Instantiations}
\subsection{Classifiers}
\[
\mathcal{L}
= \alpha\,\mathrm{CE}(y, p_\theta(\cdot \mid c;1))
+ \beta\,T^{2}\,\KL\!\big(q_\phi(\cdot \mid c;T)\,\|\,p_\theta(\cdot \mid c;T)\big).
\]

\subsection{LLMs (token-level)}
\[
\mathcal{L}
= \sum_{t}\Big[
\alpha \, \big(-\log p_\theta(y_t \mid c)\big)
+ \beta \, T^{2} \,\KL\!\big(q_\phi(\cdot \mid c;T)\,\|\,p_\theta(\cdot \mid c;T)\big)
\Big].
\]

\subsection{Protein Structure Prediction}
\[
\mathcal{L}
= \beta\,D\!\big(q_\phi(R \mid \text{seq}), p_\theta(R \mid \text{seq})\big)
+ \gamma \, \|h_T - h_S\|^2
+ \alpha \,\ell_{\text{hard}} \; \text{(e.g., RMSD/MSE to ground truth)}.
\]

\subsection{Diffusion / Score-Based Models}
\[
\mathcal{L}
= \E_{t,x_t}\!\left[
\|\epsilon_T(x_t,t,c)-\epsilon_\theta(x_t,t,c)\|^2
\right]
\quad \text{or} \quad
\E\!\left[\|s_T - s_\theta\|^2\right].
\]

\section{Proofs and Derivations}
\subsection{KL and Cross-Entropy Equivalence}
\[
\KL(q\|p) = \sum_i q_i \log \frac{q_i}{p_i} = H(q,p) - H(q).
\]
Since $H(q)$ is constant in $p$, minimizing KL is equivalent to minimizing cross-entropy.

\subsection{Temperature Scaling and $T^2$ Factor}
Let $p = \text{softmax}(z/T)$.
\[
\frac{\partial \tilde{\mathcal{L}}}{\partial z_j}
= \frac{1}{T}(p_j - q_j).
\]
Multiplying by $T^2$ rescales gradients:
\[
\frac{\partial (T^2 \tilde{\mathcal{L}})}{\partial z_j}
= T(p_j - q_j).
\]

\subsection{Sequence-Level KL Decomposition (LLMs)}
\[
\KL\big(q(y_{1:T}\mid x)\,\|\,p(y_{1:T}\mid x)\big)
= \sum_{t=1}^T
\E_{q(y_{<t}\mid x)}
\Big[\KL\big(q(y_t\mid y_{<t},x)\,\|\,p(y_t\mid y_{<t},x)\big)\Big].
\]

\subsection{MSE Minimizer is Conditional Expectation}
For $o$, $c$:
\[
g^*(c) = \arg\min_g \E[\|o-g(c)\|^2] = \E[o\mid c].
\]

\subsection{Gradient View of Distillation}
\[
\nabla_\theta \mathcal{L}_{\text{distill}}
= -\E_{c}\E_{o\sim q(\cdot\mid c)}\big[\nabla_\theta \log p_\theta(o\mid c)\big].
\]
This is maximum likelihood under teacher samples.

\section{Practical Knobs}
\begin{itemize}
    \item Balance weights $\alpha,\beta,\gamma,\delta$.
    \item Divergence choice: $\KL(q\|p)$ (mode covering) vs. $\KL(p\|q)$ (mode seeking).
    \item Temperature $T$: higher $T$ softens targets; $T^2$ scaling stabilizes gradients.
    \item Context choice $c$: teacher-forcing vs. student-rollout, diffusion schedules, etc.
\end{itemize}

\section{Pseudocode Sketches}
\subsection{LLM Token-Wise Distillation}
\begin{verbatim}
for (x, y) in data:
    with no_grad():
        q_logits = teacher(x, y_inp=y[:-1])
        q = softmax(q_logits / T)
    p_logits = student(x, y_inp=y[:-1])
    p = softmax(p_logits / T)

    loss_hard = CE_hard(y, logits=student(x, y_inp=y[:-1]))
    loss_soft = KL(q, p) * (T**2)
    loss = alpha*loss_hard + beta*loss_soft
    loss.backward(); opt.step()
\end{verbatim}

\subsection{Diffusion}
\begin{verbatim}
for x0 in data:
    t ~ Uniform({1,...,T_steps})
    eps ~ Normal(0, I)
    x_t = forward_diffuse(x0, eps, t)

    eps_T = teacher.predict_noise(x_t, t)
    eps_S = student.predict_noise(x_t, t)

    loss = mse(eps_S, eps_T)
    loss.backward(); opt.step()
\end{verbatim}

\subsection{Protein Structures}
\begin{verbatim}
for seq, gt in data:
    with no_grad():
        q_bins, h_T = teacher(seq)
    p_bins, h_S = student(seq)

    loss_kl  = KL(q_bins, p_bins)
    loss_feat = mse(h_S, h_T)
    loss_gt   = rmsd(student_coords(seq), gt_coords)
    loss = beta*loss_kl + gamma*loss_feat + alpha*loss_gt
    loss.backward(); opt.step()
\end{verbatim}

\section*{One-Line Contrast with GANs}
Teacher-student minimizes a fixed divergence to a teacher (cooperative),
while GANs define an adversarial game with a discriminator-driven divergence.
\section{Teacher-Student with RLHF (Fixed Reward Model)}
\label{sec:rlhf}

Reinforcement Learning with Human Feedback (RLHF) can be expressed in the
teacher-student framework when the reward model is fixed.

\subsection{Setup}
\begin{itemize}
    \item \textbf{Student (policy):} $\pi_\theta(y \mid x)$.
    \item \textbf{Teacher (reward-shaped distribution):}
    \[
    \pi^*(y \mid x) \propto \pi_{\text{ref}}(y \mid x)\,
    \exp\!\left(\tfrac{1}{\beta} r_\phi(x,y)\right),
    \]
    where $\pi_{\text{ref}}$ is a reference policy and $r_\phi$ is the fixed reward model.
    \item \textbf{Goal:} Distill $\pi^*$ into $\pi_\theta$ by minimizing KL divergence.
\end{itemize}

\subsection{Objective}
We maximize the reward-regularized objective:
\[
\max_{\theta} \;
\E_{x\sim\mathcal{D},\, y\sim\pi_\theta(\cdot\mid x)}
\Big[ r_\phi(x,y) - \beta \log \tfrac{\pi_\theta(y \mid x)}{\pi_{\text{ref}}(y \mid x)} \Big].
\]

Equivalently, the optimal distribution is
\[
\pi^*(y \mid x) \propto \pi_{\text{ref}}(y \mid x)\,\exp\!\left(\tfrac{1}{\beta}r_\phi(x,y)\right).
\]

The student then minimizes:
\[
\mathcal{L}_{\text{RLHF}}(\theta) = \KL\!\big(\pi^*(\cdot\mid x) \,\|\, \pi_\theta(\cdot\mid x)\big).
\]

\subsection{Gradient Derivation}
Expanding the gradient:
\[
\nabla_\theta \mathcal{L}_{\text{RLHF}}
= -\,\E_{y \sim \pi^*(\cdot\mid x)} \big[ \nabla_\theta \log \pi_\theta(y \mid x) \big].
\]

Since $\pi^*$ may be intractable to sample exactly, we use a policy-gradient estimator:
\[
\nabla_\theta J(\theta) =
\E_{y \sim \pi_\theta(\cdot\mid x)} \Big[
\big(r_\phi(x,y) - \beta \log \tfrac{\pi_\theta(y \mid x)}{\pi_{\text{ref}}(y \mid x)}\big)\,
\nabla_\theta \log \pi_\theta(y \mid x) \Big].
\]

\subsection{Interpretation}
\begin{itemize}
    \item The \textbf{teacher} is the reward-shaped distribution $\pi^*$.
    \item The \textbf{student} is $\pi_\theta$, trained to approximate $\pi^*$.
    \item The \textbf{KL term} acts as a distillation regularizer, keeping $\pi_\theta$ close to $\pi_{\text{ref}}$.
    \item This shows RLHF with a fixed reward model is a form of
    \emph{teacher-student distillation}, where human preference (via $r_\phi$)
    modifies the teacher distribution.
\end{itemize}

\subsection{Pseudocode (Fixed-Reward RLHF)}
\begin{verbatim}
# Inputs:
#   D            : supervised prompts/distribution over x
#   pi_ref       : frozen reference policy
#   r_phi        : fixed reward model (trained from human preferences)
#   pi_theta     : trainable policy (student)
#   beta         : KL regularization strength
#   eta_ent      : (optional) entropy bonus coefficient
#   b(x)         : (optional) baseline/critic to reduce variance

for epoch in 1..E:
    for minibatch x ~ D:
        # 1) Sample responses from current policy
        y ~ pi_theta(. | x)                 # e.g., nucleus/temperature sampling

        # 2) Compute rewards and logs
        r = r_phi(x, y)                     # scalar reward
        logp_theta = log pi_theta(y | x)
        logp_ref   = log pi_ref(y | x)

        # 3) Form KL-regularized advantage (reward shaping)
        #    A = r - beta * (log pi_theta - log pi_ref) - b(x)
        #    (b(x) can be moving average or learned value baseline)
        A = r - beta * (logp_theta - logp_ref) - b(x)

        # 4) Policy loss (REINFORCE-style), optional entropy bonus
        loss_policy = - mean( A * logp_theta )
        loss_entropy = - eta_ent * mean( entropy(pi_theta(. | x)) )
        loss = loss_policy + loss_entropy

        # 5) Optimize student policy
        loss.backward()
        optimizer.step(); optimizer.zero_grad()

        # (Optional) Monitor forward KL to reference:
        #   KL_forward = mean( KL(pi_theta(.|x) || pi_ref(.|x)) )
        # (Optional) Use PPO-style clipping on ratio exp(logp_theta - logp_old)
\end{verbatim}

\paragraph{Notes.}
\begin{itemize}[leftmargin=*]
    \item The KL penalty appears inside the shaped advantage
    $A = r_\phi - \beta(\log \pi_\theta - \log \pi_{\text{ref}})$,
    which yields the same stationary point as minimizing
    $\KL\!\big(\pi^* \| \pi_\theta\big)$.
    \item A baseline $b(x)$ reduces gradient variance (e.g., moving average or learned value head).
    \item Entropy bonus encourages exploration and prevents early collapse.
    \item PPO clipping, adaptive $\beta$, or trust-region updates can further stabilize training.
\end{itemize}

\end{document}
